\subsection{局限性}
\label{sec:discussion-limitations}

若干约束界定了本研究的适用范围，应指导结果解读和应用。

\textbf{训练数据需求。}该方法依赖于由低平均输入信号和高平均参考目标组成的配对训练数据。获取此类配对需要可控的采集会话，并具有足够的重复性以在同一空间位置跨多个平均级别进行配准。
该协议对于实验室试样是可行的，但为现场部署或难以进行稳定跨会话配准的材料带来了实际障碍。

\textbf{跨厚度泛化。}泛化至训练集以外厚度的铝板会引入可测量的性能下降。在 3\,mm 试样上训练的 PMDF-Net 在 5\,mm 板上显示信噪比(SNR)下降$1.8$\,dB，在 8\,mm 板上下降$3.2$\,dB，无需 retraining。这种性能下降源于 Lamb 波模态组成和群速度的厚度依赖性变化，这些变化改变了接收信号的有效中心频率。

\textbf{缺陷类型覆盖。}缺陷泛化实验使用机加工槽和边缘切口作为测试缺陷。这些几何形状产生相对较强、宽带的反射。
其他缺陷类型——特别是具有亚毫米开口位移的裂纹和复合界面的分层——可能产生较弱、更局部化的散射特征，而当前测试集未捕捉到这些特征。

\textbf{计算成本。}训练 PMDF-Net 需要 GPU 资源；物理约束损失和多域架构相对于单域基线增加了计算负担。然而，推理效率很高，在标准硬件上以商业相关速度运行。

\textbf{材料范围。}所有训练和评估均使用铝试样。碳纤维增强聚合物或其他复合层压板表现出 substantially different 各向异性波传播行为，需要单独的训练数据，
可能还需要架构 adaptation。